\documentclass[11pt,a4paper]{article}

\usepackage[T1]{fontenc}
\usepackage[utf8]{inputenc}
\usepackage[scaled=0.95]{helvet}
\renewcommand{\familydefault}{\sfdefault}
\usepackage{microtype}
\usepackage[a4paper,margin=19mm,headheight=22pt,footskip=12mm]{geometry}
\usepackage[dvipsnames,table]{xcolor}
\usepackage{hyperref}
\usepackage{bookmark}
\usepackage{enumitem}
\usepackage{titlesec}
\usepackage{fancyhdr}
\usepackage[most]{tcolorbox}
\usepackage{parskip}
\usepackage{lastpage}
\usepackage{xurl}

\definecolor{DARESNavy}{HTML}{102A43}
\definecolor{DARESTeal}{HTML}{0B7285}
\definecolor{DARESMuted}{HTML}{66788A}
\definecolor{DARESLine}{HTML}{D9E2EC}
\definecolor{DARESBlueLight}{HTML}{EAF2FF}
\definecolor{DARESTealLight}{HTML}{E7F5F7}
\definecolor{DARESOrangeLight}{HTML}{FFF4E6}
\definecolor{DARESGreenLight}{HTML}{EAF7F0}
\definecolor{DARESPaper}{HTML}{F7F9FC}
\definecolor{DARESCoral}{HTML}{E76F51}

\hypersetup{
  colorlinks=true,
  linkcolor=DARESTeal,
  urlcolor=DARESTeal,
  citecolor=DARESTeal,
  pdftitle={BLDC Motor Control - Learning and Project Documentation Guide},
  pdfauthor={DARES Electronics},
  pdfsubject={BLDC and ODrive onboarding, implementation, and documentation},
  pdfkeywords={BLDC, ODrive, ESC, six-step commutation, PWM, FOC}
}
\urlstyle{same}
\setcounter{tocdepth}{1}
\setlength{\parindent}{0pt}
\setlength{\parskip}{5pt}
\setlist[itemize]{leftmargin=5.5mm,itemsep=2.5pt,topsep=3pt}
\setlist[enumerate]{leftmargin=7mm,itemsep=3pt,topsep=3pt}
\emergencystretch=3em

\titleformat{\section}
  {\Large\bfseries\color{DARESNavy}}
  {\thesection.}{0.55em}{}
  [\vspace{2pt}\color{DARESTeal}\titlerule]
\titleformat{\subsection}
  {\large\bfseries\color{DARESTeal}}
  {\thesubsection}{0.5em}{}
\titleformat{\subsubsection}
  {\normalsize\bfseries\color{DARESNavy}}
  {\thesubsubsection}{0.5em}{}
\titlespacing*{\section}{0pt}{14pt}{8pt}
\titlespacing*{\subsection}{0pt}{10pt}{3pt}
\titlespacing*{\subsubsection}{0pt}{7pt}{2pt}

\pagestyle{fancy}
\fancyhf{}
\lhead{\small\bfseries\color{DARESTeal}DARES ELECTRONICS}
\rhead{\small\color{DARESMuted}BLDC Motor Control Documentation}
\lfoot{\scriptsize\color{DARESMuted}Learning and project documentation guide}
\rfoot{\scriptsize\color{DARESMuted}Page \thepage\ of \pageref*{LastPage}}
\renewcommand{\headrulewidth}{0.4pt}
\renewcommand{\footrulewidth}{0.4pt}
\renewcommand{\headrule}{\hbox to\headwidth{\color{DARESLine}\leaders\hrule height \headrulewidth\hfill}}
\renewcommand{\footrule}{\hbox to\headwidth{\color{DARESLine}\leaders\hrule height \footrulewidth\hfill}}

\newtcolorbox{callout}[1]{
  enhanced,
  breakable,
  colback=DARESTealLight,
  colframe=DARESLine,
  boxrule=0.4pt,
  borderline west={3pt}{0pt}{DARESTeal},
  arc=0pt,
  left=3mm,
  right=3mm,
  top=2.5mm,
  bottom=2.5mm,
  title={#1},
  coltitle=DARESNavy,
  fonttitle=\bfseries,
  attach title to upper={\par\vspace{1mm}}
}

\newtcolorbox{warningbox}[1]{
  enhanced,
  breakable,
  colback=DARESOrangeLight,
  colframe=DARESLine,
  boxrule=0.4pt,
  borderline west={3pt}{0pt}{DARESCoral},
  arc=0pt,
  left=3mm,
  right=3mm,
  top=2.5mm,
  bottom=2.5mm,
  title={#1},
  coltitle=DARESNavy,
  fonttitle=\bfseries,
  attach title to upper={\par\vspace{1mm}}
}

\newtcolorbox{resourcebox}{
  enhanced,
  breakable,
  colback=DARESBlueLight,
  colframe=DARESLine,
  boxrule=0.4pt,
  arc=0pt,
  left=3mm,
  right=3mm,
  top=2.5mm,
  bottom=2.5mm
}

\newcommand{\Resource}[3]{%
  \begin{resourcebox}
    \href{#2}{\textbf{\color{DARESTeal}#1}}\par
    {\small #3}
  \end{resourcebox}
}

\newcommand{\RepoEntry}[2]{%
  \texttt{\detokenize{#1}}\hfill{\small\color{DARESMuted}#2}\par
}

\begin{document}
\hypersetup{pageanchor=false}

\begin{titlepage}
  \thispagestyle{empty}
  \pagecolor{DARESPaper}
  \color{DARESNavy}

  \vspace*{20mm}
  {\large\bfseries\color{DARESTeal}DARES ELECTRONICS\par}
  \vspace{12mm}

  {\fontsize{30}{34}\selectfont\bfseries BLDC Motor Control\par}
  \vspace{3mm}
  {\Large\color{DARESMuted}Learning and Project Documentation Guide\par}

  \vspace{15mm}
  \begin{tcolorbox}[
    colback=DARESNavy,
    colframe=DARESNavy,
    arc=0pt,
    boxrule=0pt,
    left=4mm,right=4mm,top=4mm,bottom=4mm]
    \centering\bfseries\color{white}
    PRINCIPLE \(\rightarrow\) SIMULATION \(\rightarrow\) ODRIVE
    \(\rightarrow\) HARDWARE \(\rightarrow\) TESTING \(\rightarrow\) CONTROL
  \end{tcolorbox}

  \vspace{15mm}
  {\large\color{DARESMuted}
    A beginner-friendly reference for learning six-step ESC control,
    validating concepts on ODrive, documenting engineering work, and
    progressing toward closed-loop control and field-oriented control.
    \par}

  \vfill
  \begin{callout}{CURRENT PRIORITY}
    BLDC principle \(\rightarrow\) three-phase inverter \(\rightarrow\)
    six-step commutation \(\rightarrow\) PWM \(\rightarrow\) simulation
    \(\rightarrow\) ODrive bring-up. Use ODrive to connect the theory to
    real closed-loop control before moving to a custom ESC and advanced
    FOC development.
  \end{callout}

  \vfill
  {\color{DARESCoral}\rule{32mm}{2pt}\par}
  \vspace{3mm}
  {\small\color{DARESMuted}Version 1.1 \textbar\ September 2026\par}
\end{titlepage}
\nopagecolor
\hypersetup{pageanchor=true}

\tableofcontents
\clearpage

\section{Purpose and scope}

This guide gives new Motor Control members a shared starting point. It
explains what to learn, what to simulate, what to document, and what evidence
is required before the team moves from basic six-step commutation to
closed-loop control and FOC.

\subsection{Learning objectives}

By using this guide, a member should be able to:

\begin{itemize}
  \item Explain how a BLDC motor produces torque and why electronic commutation is required.
  \item Explain the purpose of the three-phase inverter, gate driver, PWM, dead time, and sensing circuits.
  \item Build or review a basic six-step simulation and interpret voltage, current, speed, and torque results.
  \item Bring up an ODrive safely and connect motor parameters, sensing, FOC, and closed-loop control to a real motor.
  \item Document schematic, firmware, simulation, and test work so another member can reproduce it.
  \item Recognise when the system is ready for speed feedback, current control, and FOC.
\end{itemize}

\begin{callout}{IMPORTANT BOUNDARY}
  This guide is an onboarding and documentation standard. It does not replace
  component datasheets, a design review, electrical safety procedures, or
  validation on the actual target motor and power system.
\end{callout}

\subsection{Recommended working method}

Learn enough theory to understand the physical system, build a simple model,
validate key ideas on an ODrive, compare the model with measured behaviour,
and then deepen the control theory or develop custom hardware. Theory,
simulation, hardware, and testing should not be treated as separate projects.
Each step should create evidence for the next one.

\section{Recommended learning path}

The sequence below is designed for members with no previous motor-control
experience. Complete each stage well enough to explain it in your own words
before moving to the next stage.

\subsection{BLDC motor operating principle}

A BLDC motor has a permanent-magnet rotor and three stator windings. Current
in a winding creates a stator magnetic field. The controller changes which
windings carry current so that the stator field advances around the motor.
The rotor follows this field and produces mechanical rotation. Unlike a
brushed DC motor, commutation is performed electronically by the ESC.

\textbf{Learning check}
\begin{itemize}
  \item Identify the rotor, stator, and three motor phases.
  \item Explain why current direction changes the stator magnetic field.
  \item Explain why a controller must know, or estimate, rotor position.
\end{itemize}

\Resource
  {YouTube: How does a Brushless DC Motor work?}
  {https://www.youtube.com/watch?v=Hw1q-jbhPzc}
  {Animated introduction to the rotor, stator, three-phase windings, and electronic commutation.}

\subsection{Three-phase inverter and gate driver}

The battery supplies DC power. A three-phase inverter uses three
half-bridges, normally six MOSFETs, to connect each motor phase to the
positive DC bus, the negative bus, or a floating state. The MCU generates
logic-level commands; a gate driver supplies the voltage and current needed
to switch the MOSFET gates. Dead time prevents the high-side and low-side
devices in one half-bridge from conducting simultaneously.

\textbf{Learning check}
\begin{itemize}
  \item Explain high-side, low-side, half-bridge, and shoot-through.
  \item Explain why an MCU pin should not directly drive a power MOSFET gate.
  \item Identify the roles of bootstrap circuitry, gate resistors, and dead time.
\end{itemize}

\subsection{Six-step commutation}

In six-step commutation, two phases are driven and the third phase is normally
floating. One driven phase carries current into the motor and the other
carries current out. The controller advances through six electrical states
so that the stator field moves around the motor. Reversing the order reverses
the direction. The exact switching table depends on phase naming, the Hall
sequence, and the desired direction.

\textbf{Learning check}
\begin{itemize}
  \item Describe the current path in all six states.
  \item Identify the high-side device, low-side device, and floating phase in each state.
  \item Explain why the sequence must be synchronised with rotor position.
\end{itemize}

\Resource
  {YouTube: Understanding Six-Step Commutation for BLDC Motors}
  {https://www.youtube.com/watch?v=WYJWdMV3YMs}
  {MathWorks introduction to the switching sequence and simulated behaviour of a six-step BLDC drive.}

\subsection{PWM and power control}

Pulse-width modulation switches a MOSFET rapidly. Duty cycle is the fraction
of each PWM period for which the command is active. Increasing duty cycle
generally increases the average voltage available to the motor, but current
and torque still depend on winding resistance, inductance, back-EMF, load,
commutation timing, and the selected PWM strategy. PWM frequency is a
compromise among current ripple, audible noise, switching loss, timer
resolution, and control bandwidth.

\textbf{Learning check}
\begin{itemize}
  \item Define PWM frequency, period, and duty cycle.
  \item Explain why duty cycle is not the same as guaranteed motor speed.
  \item Explain how dead time and switching losses affect real hardware.
\end{itemize}

\subsection{Rotor position and back-EMF}

Hall sensors provide coarse rotor-position sectors, while encoders provide
finer mechanical position. A sensorless six-step controller can monitor the
voltage of the floating phase and detect back-EMF zero crossings. Because
back-EMF is small near zero speed, a sensorless ESC normally needs an
open-loop starting sequence before it can transition to reliable
closed-loop commutation.

\textbf{Learning check}
\begin{itemize}
  \item Compare Hall, encoder, and sensorless position information.
  \item Explain why back-EMF increases with speed.
  \item Explain why low-speed sensorless starting is difficult.
\end{itemize}

\section{Simulation before final hardware}

Simulation is used to test understanding and reduce risk before the final
power stage is energised. It is not a substitute for measurement. The first
model should be simple enough that every block and parameter can be
explained. Suitable tools include MATLAB/Simulink, PLECS, PSIM, and LTspice.

\subsection{Minimum first simulation}

\begin{itemize}
  \item DC bus source with a documented voltage and current limit.
  \item Three-phase inverter or an equivalent switching model.
  \item BLDC motor model with documented electrical and mechanical parameters.
  \item Six-step commutation sequence and adjustable PWM duty cycle.
  \item Mechanical load or an explicitly stated no-load assumption.
\end{itemize}

\subsection{Required outputs}

\begin{itemize}
  \item Phase voltages and phase currents.
  \item DC bus voltage and current.
  \item Motor speed and electromagnetic torque.
  \item Commutation state and PWM duty cycle.
  \item Start-up transient and steady-state behaviour.
\end{itemize}

\subsection{What the report must explain}

Do not upload plots without interpretation. State the objective, assumptions,
input parameters, solver and sample time, expected behaviour, observed
behaviour, and conclusion. If the model behaves incorrectly, record the
failure and the evidence used to diagnose it. A failed simulation with a
clear explanation is more useful than an unexplained successful screenshot.

\begin{callout}{SIMULATION EXIT CRITERION}
  A member can explain the current path for each commutation state and can
  account for the main voltage, current, speed, and torque waveforms. Changing
  duty cycle or load produces a physically reasonable response.
\end{callout}

\Resource
  {Official reference: MathWorks Six Step Commutation}
  {https://www.mathworks.com/help/mcb/ref/sixstepcommutation.html}
  {Reference for six-step switching logic, Hall-state inputs, rotor-position inputs, and switching sequences.}

\section{Practical implementation with ODrive}

ODrive provides a ready-made motor-control platform containing the power
stage, gate drivers, current measurement, MCU, control firmware,
communication interface, and rotor-feedback inputs. It allows the team to
move from a simulated motor to a controlled physical motor without first
designing a complete ESC. It is useful for learning motor calibration,
current control, velocity control, position feedback, FOC, regenerative
braking, and fault diagnosis.

\begin{callout}{IMPORTANT DISTINCTION}
  ODrive is not a six-step commutation teaching board. Its normal control
  approach is closed-loop current control and FOC. Six-step commutation
  should still be learned and simulated because it makes the inverter,
  phase-current paths, PWM, and rotor-position problem easier to understand.
  ODrive is the next practical step, not a replacement for those fundamentals.
\end{callout}

\subsection{How ODrive maps to the motor-control system}

\begin{itemize}
  \item DC input and bus capacitors provide energy to the inverter.
  \item Gate drivers and MOSFET bridges switch the three motor phases.
  \item Current shunts and amplifiers provide phase-current feedback for the inner current loop.
  \item The MCU executes calibration, FOC, protection, and torque, velocity, or position control.
  \item An encoder, Hall sensors, or a sensorless estimator provides rotor information.
  \item USB, UART, CAN, or other interfaces allow configuration, commands, logging, and fault reporting.
  \item The brake resistor or another regeneration strategy manages energy returned to the DC bus.
\end{itemize}

\subsection{Record the exact hardware and firmware first}

Before following any tutorial, record the ODrive model, voltage version, board
revision, firmware version, host-tool version, and motor-feedback method.
ODrive v3.x firmware 0.5.x and current ODrive products running 0.6.x do not
use identical configuration names or procedures. Treat a video as a
conceptual demonstration and use documentation that matches the exact board
and firmware for actual commands.

\subsection{Recommended bring-up workflow}

\begin{enumerate}
  \item Inspect the board and wiring. Secure the motor, remove the propeller, and confirm DC supply polarity and voltage range.
  \item Connect the three motor phases and the selected encoder or Hall feedback. Confirm pole pairs and feedback resolution.
  \item Set conservative bus-voltage, motor-current, calibration-current, velocity, and acceleration limits before enabling the axis.
  \item Plan for regenerated energy. Configure the correct brake resistor or regeneration limit for the board and power source.
  \item Run motor calibration so the controller can estimate phase resistance and inductance. Investigate errors rather than repeatedly increasing limits.
  \item Calibrate and verify the encoder or configure the intended sensorless method. Confirm that position and velocity change in the expected direction.
  \item Begin closed-loop operation at low torque and speed. Test start, stop, direction, and loss-of-feedback behaviour before increasing the command.
  \item Log DC bus voltage, current, measured speed, target speed, controller mode, temperature, and error history.
\end{enumerate}

\subsection{What members should learn from ODrive}

\begin{itemize}
  \item Why motor resistance and inductance are needed by a current controller.
  \item How pole pairs relate mechanical rotation to electrical angle.
  \item How encoder alignment and direction affect torque production.
  \item How torque, velocity, and position modes use nested control loops.
  \item Why current, velocity, bus-voltage, temperature, and fault limits are part of the controller design.
  \item Why deceleration can return energy to the DC bus and cause over-voltage or regeneration faults.
\end{itemize}

\subsection{Required ODrive evidence}

\begin{itemize}
  \item Board model, firmware version, power source, motor model, pole pairs, and feedback device.
  \item Saved configuration or a script that reproduces every changed parameter.
  \item Calibration results, including measured resistance and inductance where available.
  \item A low-speed closed-loop test with target and measured speed, bus voltage, current, and temperature.
  \item Complete error output and the reasoning used to diagnose each fault.
  \item A conclusion explaining what the ODrive test proves and what it does not prove about the custom ESC.
\end{itemize}

\Resource
  {YouTube: Unpacking and testing the new product ODrive}
  {https://www.youtube.com/watch?v=DVN25PbMfDE\&list=PLc2RScfrSFEDcQp6Qvc7mBNre8ABCoz9p}
  {Practical introduction to connecting, configuring, and testing a BLDC motor with an ODrive controller.}

\Resource
  {Official ODrive reference: Motor Calibration Parameters}
  {https://docs.odriverobotics.com/v/latest/articles/motor-parameters.html}
  {Explains calibration current and voltage, resistance and inductance measurement, and common calibration faults.}

\Resource
  {Official ODrive reference: Hardware Configuration}
  {https://docs.odriverobotics.com/v/latest/manual/hardware-config.html}
  {Covers power limits, regenerated energy, braking, temperature monitoring, encoders, and sensorless configuration.}

\begin{warningbox}{ODRIVE EXIT CRITERION}
  The configuration is reproducible; calibration completes without an
  unexplained fault; the unloaded motor starts, runs, changes speed, and stops
  predictably at conservative limits; measured values are logged; and the
  member can relate board behaviour to inverter, sensing, and control concepts.
\end{warningbox}

\section{ESC hardware documentation}

Hardware design should begin after the team understands the power path and
has a basic model. The schematic must be divided into functional blocks, and
every major component must have a documented selection reason. Copying a
reference design is acceptable only when the team verifies that its voltage,
current, thermal, and control assumptions match the project.

\subsection{Required hardware blocks}

\subsubsection{MCU and timing}
Document PWM channels, complementary outputs, dead-time support, ADC
triggering, interrupt timing, communication, and available processing margin.

\subsubsection{Gate driver}
Document high-side and low-side drive capability, bootstrap limits, UVLO,
fault reporting, propagation delay, peak gate current, and supported
switching frequency.

\subsubsection{MOSFET power stage}
Document drain-source voltage margin, continuous and pulse current ratings,
\(R_{\mathrm{DS(on)}}\), gate charge, switching loss, thermal resistance,
package, and PCB current paths.

\subsubsection{Current sensing}
Document shunt location, resistance, power rating, amplifier range and
bandwidth, common-mode limits, ADC range, filtering, and over-current
protection.

\subsubsection{Voltage sensing}
Document battery and DC bus dividers, phase-voltage or back-EMF scaling,
filtering, transient protection, and ADC input limits.

\subsubsection{Power supplies}
Document logic rails, gate-driver supply, regulator input range, start-up
behaviour, decoupling, and power sequencing.

\subsubsection{Protection and thermal design}
Document reverse polarity, input transients, shoot-through prevention,
over-current, under-voltage, over-voltage, over-temperature, fusing, and heat
removal.

\subsubsection{Interfaces}
Document rotor-sensor inputs, command input, telemetry, programming, debug,
and fault/status outputs.

\subsection{Component selection record}

For each critical component, record the selected part, required rating,
calculated margin, relevant datasheet pages, alternatives considered, reason
for selection, risks, and validation plan. Absolute maximum ratings are not
normal operating targets; design against recommended operating conditions and
expected transients.

\section{Firmware documentation}

Firmware documentation must describe the control timing and the relationship
between software commands and physical switching. A reader should be able to
trace a target command through commutation logic, PWM generation, ADC
sampling, feedback calculation, protection checks, and fault response.

\subsection{Minimum topics}

\begin{itemize}
  \item Timer, PWM, complementary output, dead-time, and update-event configuration.
  \item Six-step state machine, direction control, and commutation timing.
  \item Hall, encoder, or back-EMF input processing and validation.
  \item ADC trigger timing, scaling, offset calibration, and filtering.
  \item RPM calculation and timeout behaviour.
  \item Start-up, run, stop, and fault states.
  \item Over-current, over-voltage, under-voltage, and over-temperature responses.
  \item Communication protocol, units, limits, and error handling.
\end{itemize}

\subsection{Function-level documentation}

Important functions should state their purpose, inputs, outputs, units, valid
range, call rate, side effects, and safety assumptions. Constants such as pole
pairs, current-sensor gain, ADC reference, timer clock, and PWM frequency must
have one authoritative definition. Changes to these constants should be
reviewed because they can alter physical limits.

\begin{callout}{FIRMWARE EXIT CRITERION}
  The motor can be commanded to start and stop predictably; commutation and
  feedback units are verified; every defined fault forces a documented safe
  state; and the firmware version is recorded with each test.
\end{callout}

\section{Bench testing and safety}

The first custom-hardware objective is deliberately narrow: rotate the motor
reliably using six-step commutation and PWM. Test one assumption at a time and
capture the evidence. Start with logic-level and gate-signal checks before
connecting the motor to the complete power stage.

\subsection{Minimum safety controls}

\begin{itemize}
  \item Remove the propeller unless an approved guarded thrust test specifically requires it.
  \item Use a current-limited supply where practical, plus an appropriate fuse and accessible emergency disconnect.
  \item Secure the motor and loose wiring before energising the system.
  \item Verify PWM polarity, complementary outputs, dead time, and gate waveforms before applying full bus voltage.
  \item Use differential or isolated measurements where required; do not attach an earth-referenced oscilloscope ground to an unsafe switching node.
  \item Begin at low voltage and conservative current limits, then increase only after reviewing measurements.
  \item Monitor MOSFET, gate driver, motor, connector, and cable temperature.
  \item Stop immediately for unexpected current, noise, heating, loss of commutation, or unstable control.
\end{itemize}

\subsection{Test record}

\begin{itemize}
  \item Test ID, date, operator, reviewer, and objective.
  \item Schematic/PCB revision, firmware commit, configuration, and instrument list.
  \item Power supply voltage and current limit, motor, sensor, load, and propeller status.
  \item Procedure, expected result, measured result, plots or oscilloscope captures, and pass/fail decision.
  \item Unexpected behaviour, likely cause, follow-up action, and linked issue or design note.
\end{itemize}

\begin{warningbox}{FIRST HARDWARE MILESTONE}
  The motor starts, continues rotating, and stops on command. Current remains
  within the planned limit, temperatures remain acceptable, gate signals show
  no shoot-through, and the result is reproducible.
\end{warningbox}

\section{Closed-loop speed and current control}

\subsection{Speed feedback}

Open-loop control applies a command but does not guarantee the requested
speed. Closed-loop speed control compares target RPM with measured RPM and
adjusts the command to reduce the error. Documentation must state how RPM is
measured, the update rate, filtering, units, valid range, and behaviour when
feedback is lost.

\subsection{PID understanding}

Proportional action reacts to present error. Integral action accumulates error
and can remove steady-state offset, but it needs output limiting and
anti-windup. Derivative action reacts to rate of change and is sensitive to
measurement noise. A PI controller is often sufficient for speed control; the
correct choice depends on the plant, sample rate, noise, and performance
requirement. Gains must be recorded with the test conditions and must not be
treated as universal constants.

\subsection{Current control}

Motor torque is strongly related to current, so a current loop provides more
direct torque control and protects the power stage. In a cascaded controller,
the speed loop normally generates a current or torque request, while the
faster inner current loop controls the electrical response. The current
sensing topology and PWM/ADC timing determine which measurements are valid.

\subsection{Required control evidence}

\begin{itemize}
  \item Target, measured value, error, controller output, saturation limits, and sample time.
  \item Step response under at least two loads or operating points.
  \item Rise time, overshoot, settling behaviour, steady-state error, and oscillation/noise observations.
  \item Controller gains, tuning method, anti-windup method, filters, and safety limits.
\end{itemize}

\section{Introduction to field-oriented control}

FOC should begin only after phase-current measurement, rotor angle, PWM
timing, and basic feedback control are understood. FOC transforms measured
three-phase currents into a rotating reference frame aligned with the rotor.
This produces d-axis and q-axis current components controlled by fast PI
loops. In simplified terms, q-axis current is mainly associated with torque,
while d-axis current is associated with flux. Commanded voltages are
transformed back and applied through a modulation method such as SVPWM.

\subsection{Concepts to learn}

\begin{itemize}
  \item Mechanical angle, electrical angle, and pole pairs.
  \item Phase-current reconstruction and current-sampling windows.
  \item Clarke and Park transforms, plus their inverse transforms.
  \item d-axis and q-axis current references and PI current controllers.
  \item Voltage limiting, decoupling, and PI anti-windup.
  \item SVPWM, encoder alignment, and sensorless observers.
\end{itemize}

\Resource
  {YouTube: Basics of Field-Oriented Control}
  {https://www.youtube.com/watch?v=j0XfWufMM7s}
  {STMicroelectronics introduction to BLDC/PMSM control, FOC, current sensing, and sensor/sensorless topologies.}

\Resource
  {YouTube: Field-Oriented Control on STM32 From Scratch}
  {https://www.youtube.com/watch?v=AmS22zwl2EA}
  {Practical reference covering encoder feedback, current measurement, transforms, and SVPWM.}

\Resource
  {Official course: STM32 Motor Control Part 5}
  {https://www.st.com/content/st_com/en/support/learning/stm32-moocs/Motor-Control-Part-5-STM32-Field-Oriented-motor-control-training.html}
  {ST training covering FOC theory, STM32 implementation, current sensing, observers, and practical tools.}

\begin{callout}{FOC ENTRY CRITERION}
  The team has trustworthy current measurements, a verified electrical angle,
  deterministic PWM and ADC timing, safe current limits, and a working test
  method. FOC should not be used to hide unverified hardware behaviour.
\end{callout}

\section{Repository structure}

The repository should make the current design state easy to find. The main
README is an entry point; detailed technical work belongs in focused Markdown
documents and version-controlled project files.

\begin{tcolorbox}[
  breakable,
  colback=DARESPaper,
  colframe=DARESLine,
  boxrule=0.4pt,
  arc=0pt,
  left=3mm,right=3mm,top=3mm,bottom=3mm]
\RepoEntry{README.md}{project overview, status, and quick start}
\RepoEntry{docs/motor-basics/}{learning notes and reviewed explanations}
\RepoEntry{docs/control-theory/}{speed, current, FOC, and tuning notes}
\RepoEntry{docs/datasheets/}{links or approved copies of key datasheets}
\RepoEntry{docs/design-notes/}{engineering decisions and trade-offs}
\RepoEntry{simulation/}{model files, parameters, and reports}
\RepoEntry{hardware/schematic/}{schematic source, exports, and review notes}
\RepoEntry{hardware/pcb/}{layout source, manufacturing files, and reviews}
\RepoEntry{firmware/}{source code, configuration, and build instructions}
\RepoEntry{testing/odrive/}{ODrive configuration, scripts, logs, and faults}
\RepoEntry{testing/}{test plans, raw data, logs, and reports}
\RepoEntry{results/}{approved summary plots and milestone evidence}
\end{tcolorbox}

\subsection{File and commit expectations}

\begin{itemize}
  \item Use descriptive filenames and include units or dates where they prevent ambiguity.
  \item Do not upload an unexplained screenshot, model, binary, or data file.
  \item Link research sources and identify the exact datasheet revision used.
  \item Record failures as well as successes; do not remove evidence because a test failed.
  \item Keep generated exports separate from editable source files.
  \item Use issues or design notes for unresolved questions and link the relevant commit or test report.
  \item Request review for changes that affect switching, sensing, current limits, protection, or safety.
\end{itemize}

\section{Required document templates}

\subsection{Simulation note}
\begin{itemize}
  \item Objective and question being tested.
  \item Model scope, assumptions, parameter sources, solver, and sample time.
  \item Inputs and expected behaviour.
  \item Results with units and explanation.
  \item Limitations, problems, conclusion, and next action.
\end{itemize}

\subsection{Design decision record}
\begin{itemize}
  \item Decision and date.
  \item Requirement or problem being addressed.
  \item Alternatives considered and evidence reviewed.
  \item Selected option, calculations, margins, benefits, and disadvantages.
  \item Risks, validation plan, owner, and review status.
\end{itemize}

\subsection{Test report}
\begin{itemize}
  \item Test ID, objective, operator, reviewer, date, and location.
  \item Hardware revision, firmware commit, configuration, motor, load, supply, and instruments.
  \item Safety controls and pre-test checks.
  \item Procedure, expected result, measured data, and uncertainty where relevant.
  \item Pass/fail decision, fault observations, conclusion, and linked follow-up issue.
\end{itemize}

\subsection{Datasheet review}

Before using a component, review at least the absolute maximum ratings,
recommended operating conditions, pin descriptions, electrical
characteristics, timing requirements, thermal information, and typical
application circuit. Record which limits are design requirements and which
must be validated during testing.

\section{First tasks for new members}

These tasks turn passive learning into reviewable engineering output. New
members should work with an experienced member, but they should still form and
test their own explanation.

\subsection{Task 1 - Principle note}
Write a short explanation of how a BLDC motor, inverter, six-step sequence,
and PWM work together. Include the current path for each commutation state in
words or a switching table.

\subsection{Task 2 - Simulation}
Create or reproduce a basic six-step simulation. Change duty cycle and load,
then explain the effect on current, speed, and torque.

\subsection{Task 3 - ODrive bring-up}
Record the exact board and firmware, create a safe configuration, complete
calibration, run a conservative closed-loop speed test, and document all
measured values and faults.

\subsection{Task 4 - Hardware block review}
Choose one block -- gate driver, MOSFET stage, current sensing, voltage
sensing, power supply, or protection -- and document its requirements and one
candidate implementation.

\subsection{Task 5 - GitHub submission}
Upload the editable source, readable result, ODrive configuration or script,
and a Markdown explanation. Link every source and state what remains
uncertain.

\subsection{Task 6 - Short presentation}
Present what was learned, one mistake or unexpected result, how it was
investigated, and the next test that would reduce uncertainty.

\begin{callout}{REVIEW STANDARD}
  The goal is not to appear correct on the first attempt. The goal is to make
  assumptions visible, test them, record the evidence, and improve the next
  design.
\end{callout}

\section{Milestones and definition of done}

\subsection{Milestone 1 - Fundamentals}
A member can explain the motor, inverter, PWM, six-step current paths, and
position feedback without copying a diagram or script.

\subsection{Milestone 2 - Simulation}
The model runs, inputs and parameters are documented, waveforms are physically
reasonable, and results are explained.

\subsection{Milestone 3 - ODrive reference-platform bring-up}
The ODrive configuration is reproducible; calibration is understood; the
motor runs under conservative closed-loop limits; data and faults are
documented.

\subsection{Milestone 4 - Open-loop or commutated custom bench run}
The motor starts, runs, and stops safely; current and temperature stay within
limits; gate signals and commutation are verified; the test is reproducible.

\subsection{Milestone 5 - Closed-loop speed}
RPM feedback is validated; the controller tracks commands over the agreed
operating range; gains, limits, and response measurements are documented.

\subsection{Milestone 6 - Current control and FOC readiness}
Current sensing and electrical angle are verified; PWM/ADC timing is
deterministic; protection is tested; and the team has an agreed FOC
validation plan.

\subsection{Project-level definition of done}
\begin{itemize}
  \item Editable design sources and build instructions are committed.
  \item Requirements, calculations, assumptions, and safety limits are documented.
  \item Relevant simulation and test evidence is stored with units and version information.
  \item Known limitations and unresolved risks are explicit.
  \item Another team member can reproduce the result without relying on the original author.
  \item A reviewer has approved the milestone against its stated exit criteria.
\end{itemize}

\section{Learning resources}

Use these resources as starting points, then confirm implementation details
against the relevant motor, MCU, gate-driver, MOSFET, sensor, and regulator
datasheets. Video explanations are useful for intuition but are not the final
authority for component limits or safety decisions.

\Resource
  {YouTube - How does a Brushless DC Motor work?}
  {https://www.youtube.com/watch?v=Hw1q-jbhPzc}
  {Beginner animation for motor construction and electronic commutation.}

\Resource
  {YouTube - Understanding Six-Step Commutation for BLDC Motors}
  {https://www.youtube.com/watch?v=WYJWdMV3YMs}
  {Six-step switching and simulated motor behaviour from MathWorks.}

\Resource
  {MathWorks - Six Step Commutation reference}
  {https://www.mathworks.com/help/mcb/ref/sixstepcommutation.html}
  {Official switching-sequence reference for Hall or position inputs.}

\Resource
  {YouTube - Unpacking and testing the new product ODrive}
  {https://www.youtube.com/watch?v=DVN25PbMfDE\&list=PLc2RScfrSFEDcQp6Qvc7mBNre8ABCoz9p}
  {Practical ODrive connection, configuration, calibration, and motor test introduction.}

\Resource
  {YouTube - Basics of Field-Oriented Control}
  {https://www.youtube.com/watch?v=j0XfWufMM7s}
  {FOC introduction from STMicroelectronics.}

\Resource
  {STMicroelectronics - Motor Control Part 5}
  {https://www.st.com/content/st_com/en/support/learning/stm32-moocs/Motor-Control-Part-5-STM32-Field-Oriented-motor-control-training.html}
  {Structured course covering FOC theory, STM32 implementation, sensing, and observers.}

\Resource
  {YouTube - Field-Oriented Control on STM32 From Scratch}
  {https://www.youtube.com/watch?v=AmS22zwl2EA}
  {Practical implementation reference for current sensing, encoder feedback, transforms, and SVPWM.}

\subsection{Useful search terms}
\begin{itemize}
  \item BLDC motor working principle
  \item three-phase inverter and half-bridge explained
  \item MOSFET gate driver dead time and shoot-through
  \item BLDC six-step commutation
  \item BLDC PWM strategies
  \item back-EMF zero-crossing sensorless BLDC
  \item ODrive motor calibration, encoder setup, regeneration, and brake resistor
  \item BLDC speed PI control
  \item Clarke Park transform motor control
  \item FOC current sensing and SVPWM
\end{itemize}

\subsection{Start this week}
\begin{itemize}
  \item Watch the BLDC principle, six-step commutation, and ODrive bring-up videos.
  \item Write a one-page explanation in your own words.
  \item Build or review a basic six-step simulation and change one parameter.
  \item Create a version-specific ODrive safety checklist, then plan one conservative closed-loop test.
  \item Upload editable sources and evidence, then present one unresolved question and the next test.
\end{itemize}

\end{document}
